\section{Conclusion}
\label{sec:conclusion}

Reviewing a BRSR report means following a claim to the evidence behind it, and then to a summary for the whole company. GRASP supports this in one workflow. Each report is stored in a structured form, and the principle labels come from the structure of the report and not from a person reading the text. A classifier sorts text segments into the nine principles, a detector flags vague wording, and numerical checks create review records. A four-part Transparency Score then summarises each company. Text segments and facts remain tied to the same report, so an analyst can always return to the source. The source of every label stays visible during evaluation, and every score component has an explicit formula. The vagueness detector built for the flagging step is used again to calculate Clarity, which carries a weight of 20\% in the score. The numerical review records stay outside the score as separate outputs.

The classifier was tested on 6,201 text segments from 175 companies that were not used for supervised training. Here, 6,201 is a count of text segments, and the test group does not contain 6,201 reports or 6,201 companies. For 93.00\% of these text segments, the classifier predicts the same principle as the weak label (95\% confidence interval: 92.04--93.89\%). A confidence interval is a range that is likely to contain the true value. The macro-F1 is 0.9158, which is the average F1 score over the nine principles, so a small principle counts as much as a large one. The result depends on where the label came from. Agreement is 63.35\% for explicitly tagged narratives, 95.15\% for concept-mapped narratives, and 90.72\% for principle notes. A weak label is taken from the structure of the report, so each of these values measures agreement with the recorded assignment. None of them is an independent judgement of what the text means.

The other outputs were tested separately. On 80 development examples, the fixed vagueness detector reaches an F1 score of 0.7436. The Transparency Score was compared with 20 historical human ratings, and the Pearson correlation is $r=0.6314$ (95\% confidence interval: 0.2683--0.8519). This correlation shows how closely two sets of scores rise and fall together. The conditions behind each result are described in Section~\ref{sec:evaluation}. In short, the results cover principle assignment, language screening, and score comparison, and each one should be read only for the task it measures.

The saved outputs cover 1,164 companies in 72 recorded sectors and are presented at four levels: one worked company, selected companies within a sector, the complete financial-sector group, and all sectors. Financial-sector scores range from 29.83 to 67.56. Among the 47 sectors with at least five companies, Wearing Apparel has the lowest mean score, 42.33, and Warehousing and Transport Support has the highest, 54.55. Eight companies have reports for two years, and their mean score moved from 45.17 to 50.40, a gain of 5.23 points. Individually, six of the eight companies gained and two lost. For each company, the component breakdown names the part of the disclosure that drove the change. For example, Completeness explains most of the gain for ESAF Small Finance Bank, which rises by 11.51 points. Together, the company, sector, and year views allow an analyst to inspect the results both in detail and in summary.

The work can be extended in several directions. The first concerns the data. Collection can be widened to include annexures, charts, and additional annual reports. The second concerns the models. New comparisons need larger datasets with principle labels. Blind human-rating studies would also help, because reviewers would rate reports without seeing the automated scores. The language adaptation step can be repeated with text from declared training companies only, and then tested on newly collected companies. Business descriptions are a further source of information. A description of what a company does can make the sector assignment more detailed, and the same description can show which reporting items apply to the company.

The numerical checks have room to grow as well. At present, an analyst compares the definition of the measure, the unit, the period, and the business boundary by hand. A future check could make these comparisons before a review record is created. More annual reports would also help. With several years available, the same companies can be compared again and again under one recorded score version. Finally, the report-based approach can be tested on international sustainability standards. For each standard, the reporting items must first be mapped one by one, and the results must then be evaluated separately.
